\documentclass[journal]{IEEEtran}

\usepackage[utf8]{inputenc}
\usepackage[T1]{fontenc}
\usepackage{amsmath,amssymb,amsfonts}
\usepackage{graphicx}
\usepackage{booktabs}
\usepackage{array}
\usepackage{cite}
\usepackage{url}

\title{Supplementary Material for\\
       ``Content-Bound Verified Peer Records for Adversarially Robust
       Battery-Safety Agents in Industrial IoT''}

\author{%
\IEEEauthorblockN{Mohamed Abubakar Dini Rage, Muhammad Sannan Khaliq,
                  Md Mehedi Hasan Somrat, \\
                  Subroto Kumar Ghosh, Dong-Seong Kim, Jae Min Lee}%
}

\begin{document}
\maketitle

\section{Provenance-Substrate Platform Comparison}
\label{sec:supp-platforms}

Table~\ref{tab:blockchain-comparison} evaluates seven permissioned and
consortium platforms against the three substrate requirements of the
main paper (R1 deterministic finality, R2 zero marginal logging cost,
R3 native EVM\,/\,Solidity with receipt-query verification); PureChain
is the only one meeting all of R1--R3 with a live IIoT-scale mainnet.
Quorum and Besu expose the same
\texttt{eth\_getTransactionReceipt} interface but only \emph{configure}
gas to zero (R2) and offered no live IIoT deployment for latency
measurement (``not meas.''); Fabric, Corda, and Sawtooth reach
integrity through non-EVM APIs (gRPC, Flow, REST) that fail R3.
VeChain is EVM-native but meters VeThor per transaction (R2) and is
slow: its Thor REST receipt query is 45$\times$ slower than
PureChain's, so the five per-decision hash checks of the main paper
would add ${\approx}$4.6\,s, impractical for the online pipeline.

Receipt-query latencies for PureChain and VeChain were measured in
this work against their live mainnets.  A warm \texttt{web3.py}
session of $n=60$ queries to committed PureChain hashes gave mean
$22.7$\,ms, median $20.1$\,ms, and $p_{95}$ of $39$\,ms with no
errors, and the same campaign committed $570$ events in $9.0$\,s
($2{,}082$ total on-chain events for this work, including the 337
peer-anchoring transactions of the content-bound condition, no
failures); $n=25$ VeChain Thor REST receipt queries gave a median of
$913.9$\,ms ($p_{95}$ $1{,}228$\,ms).  Throughput and commit-latency
ranges for the remaining platforms are taken from published
surveys~\cite{belotti2019vademecum,liu2019performance} and vary with
block size, endorsement policy, and cluster scale.

A few provenance notes apply to the table.  PureChain's throughput is
the platform benchmark of ${\approx}2{,}336$ transactions per second
(TPS) sustained and ${\approx}7{,}000$\,TPS
peak~\cite{kim2026purechain}; the $570$ events committed in $9.0$\,s
above reflect our sporadic logging workload, not a capacity limit.
VeChain's ${\approx}100$--$165$\,TPS reflects 2024 PoA\,2.0 mainnet
statistics (launched 2021) that the 2019 surveys predate.
Hyperledger Fabric's BFT (SmartBFT) ordering service appears only in
v3.0~(2024), whereas v2.x supports Raft (crash-fault-tolerant)
consensus alone, and Hyperledger Sawtooth reaches deterministic
finality only in PBFT mode (PoET uses probabilistic leader election at
${\approx}3$--$5$\,TPS) and provides EVM support only through the
non-native Seth extension, so it satisfies R3 only partially.

\begin{table*}[!tbp]
\caption{Comparison of permissioned and consortium blockchain platforms
for inference-time record verification.  Receipt-query latencies for
PureChain and VeChain are measured in this work; throughput and
commit-latency ranges are from published
surveys~\cite{belotti2019vademecum,liu2019performance}.}
\label{tab:blockchain-comparison}
\centering
\resizebox{\linewidth}{!}{%
\begin{tabular}{llllllll}
\toprule
\textbf{Property}
  & \textbf{PureChain}~\cite{kim2025purechain,kim2026purechain}
  & \textbf{H.\,Fabric}
  & \textbf{Quorum}
  & \textbf{H.\,Besu}
  & \textbf{R3 Corda}
  & \textbf{VeChain}
  & \textbf{H.\,Sawtooth} \\
\midrule
Consensus
  & PoA\textsuperscript{2} & Raft\,/\,SmartBFT & IBFT\,2.0\,/\,QBFT
  & IBFT\,2.0\,/\,QBFT & Notary (BFT) & PoA\,2.0 & PoET\,/\,PBFT \\
Transaction finality
  & Deterministic & Deterministic & Deterministic & Deterministic
  & Deterministic & Deterministic & Deterministic (PBFT) \\
Write throughput (TPS)
  & ${\approx}$2{,}336 (7{,}000 peak) & 1{,}000--3{,}500 & 150--300 & 200--500
  & 100--170 & ${\approx}$100--165 & 30--100 (PBFT) \\
Commit latency (ms)
  & ${\approx}$2{,}000 & 500--2{,}000 & 1{,}000--2{,}000 & 1{,}000--2{,}000
  & 1{,}000--3{,}000 & ${\approx}$10{,}000--30{,}000 & 1{,}000--3{,}000 \\
Receipt-query latency
  & \textbf{20.1\,ms} {\scriptsize(p95\,39)} & N/A (gRPC) & not meas.
  & not meas. & N/A (Flow) & 913.9\,ms {\scriptsize(p95\,1{,}228)}
  & N/A (REST) \\
Fee model
  & \textbf{Zero} (design) & Zero & Config.\,(0) & Config.\,(0) & Zero
  & VeThor energy & Zero \\
Smart contract
  & Solidity (EVM) & Go\,/\,Java\,/\,JS & Solidity (EVM) & Solidity (EVM)
  & Kotlin\,/\,Java & Solidity (EVM) & Python\,/\,Go\,/\,JS \\
Privacy
  & PoA whitelist & Channels, private data & Tessera & Tessera\,/\,Orion
  & Need-to-know & Limited & None native \\
Live IIoT mainnet
  & \textbf{Yes} {\scriptsize(2{,}082\,txs)} & No & No & No & No & Public
  & No \\
\midrule
R1 -- Deterministic finality
  & \textbf{Yes} & Yes & Yes & Yes & Yes & Yes & Yes \\
R2 -- Zero marginal logging cost
  & \textbf{Yes} & Yes & Config. & Config. & Yes & \textbf{No} & Yes \\
R3 -- Native EVM\,/\,Solidity
  & \textbf{Yes} & \textbf{No} & Yes & Yes & \textbf{No} & Yes & Partial \\
\midrule
\textbf{All R1--R3 + live mainnet}
  & \textbf{All met} & Partial (R3) & No mainnet & No mainnet
  & Partial (R3) & Partial (R2) & Partial (R3) \\
\bottomrule
\end{tabular}%
}
\end{table*}

\section{Legacy Supply-Chain Thermal Checker}
\label{sec:supp-checker}

For the legacy Physical-Internet supply-chain task, the deterministic
physics checker evaluates whether a carrier and route choice remains
within the allowable thermal envelope.  The safe arrival temperature
threshold is computed as
\begin{equation}
\label{eq:t-safe}
T_{\mathrm{safe}}
=
45 - 0.2\max(0,80-d),
\end{equation}
and the estimated arrival temperature is
\begin{equation}
\label{eq:t-arr}
T_{\mathrm{arr}}
=
T_{\mathrm{amb,max}}
+
V_{\kappa}
+
0.1\max(0,80-d),
\end{equation}
where \(d\) is residual battery capacity in percent,
\(T_{\mathrm{amb,max}}\) is the maximum ambient temperature along the
route, and \(V_{\kappa}\) is the carrier-class thermal-variance term.
A route and carrier assignment is thermally feasible when
\begin{equation}
\label{eq:thermal-feasible}
T_{\mathrm{arr}} \leq T_{\mathrm{safe}}.
\end{equation}
The thermal thresholds are anchored to the literature on lithium-ion
thermal abuse and battery-safety
standards~\cite{iec62133,feng2018thermal,finegan2017characterising}.

\bibliographystyle{IEEEtran}
\bibliography{references}

\end{document}
